\section{Machine-learning methodology}\label{sec:methods}

\subsection{Estimator}
Gradient-boosted regression trees (XGBoost~\cite{xgboost2016}) are used throughout. The estimator handles
missing values natively, which matters here because several descriptors are reported by a
minority of papers, and it is robust on the sample sizes involved. Categorical descriptors
are one-hot encoded with an encoder fitted on the training partition only and
\texttt{handle\_unknown="ignore"}, so a level appearing solely in a held-out paper cannot
influence the encoding.

\subsection{Paper-grouped cross-validation}
Ordinary random splitting is inadmissible here. Multiple rows originate from the same
publication and share its cell, protocol and often its exact descriptor values, so a random
split would place near-duplicates on both sides of the partition and report an optimistic
score. All evaluation therefore uses \texttt{GroupKFold} on \texttt{paper\_id} with five
folds: rows from one publication are never split across training and test.

\subsection{Hyper-parameter selection}
Hyper-parameters are chosen by randomised search inside an \emph{inner} paper-grouped
split of the training partition of each outer fold, implemented with
scikit-learn~\cite{sklearn2011}. The outer test papers never influence
model selection, encoding, or any preprocessing statistic. A fixed seed of 42 is used for
NumPy, Python \texttt{random}, XGBoost, the cross-validation splitters and the synthetic
generator.

\subsection{Two models with different purposes}
This is the central methodological distinction of the study and it is maintained throughout.

\begin{itemize}
\item The \textbf{paper-grouped cross-validation} above estimates \emph{predictive
transferability} to publications the model has not seen. It is the only number quoted as a
performance estimate.
\item A \textbf{final real-only interpretation model} is fitted to all 118 usable
observations and is used \emph{solely} to characterise multivariate descriptor
associations within the compiled corpus. Its in-sample fit
($R^2 = 0.624$) is reported here only for
completeness and is not a generalisation estimate.
\end{itemize}

Conflating the two would be the most consequential error available in this analysis, so the
two models are named, stored and reported separately
(\texttt{results\_final/models/model\_card.md}).

\subsection{Metrics}
$R^2$, MAE and RMSE are reported per fold and aggregated as mean, SD and median. Because
$R^2$ is normalised by the variance of each fold's own held-out targets, folds whose papers
happen to agree with one another are scored against a very small denominator; pooled metrics
computed over all held-out predictions jointly are therefore also reported.
